\section{Current Status}
\label{sec:CurrentStatus}

\subsection{Maturity}

Active prototype with real executable output. The generator is capable and
regression-covered; generated projects are usable for local exploration.

\subsection{Solid}

\begin{itemize}
\item generator refactored from monolith to multiple templates and helper classes,
\item dry-run and compile-level tests cover representative BPMN features,
\item core runtime contracts are unit-tested,
\item monitoring topology with global-table query model,
\item interrupting boundary cancellation via in-memory registry,
\item GitHub Actions CI on push/PR (JDK 21),
\item \texttt{SendTask}, \texttt{ReceiveTask}, \texttt{ScriptTask}, \texttt{BusinessRuleTask} generate dedicated handlers,
\item inclusive gateway (OR) split/join,
\item multi-instance loop characteristics detected and reported,
\item 6 data pipeline plugins with interface-based dispatch.
\end{itemize}

\subsection{Limits}

\begin{itemize}
\item BPMN subset (no complex gateways, no data objects),
\item advanced scope semantics approximated,
\item event subprocess semantics scoped, not generalized to all top-level cases,
\item history retention is intentionally coarse.
\end{itemize}

\subsection{Open decisions}

\begin{description}
\item[Top-level event subprocesses] Whether timer, error, and escalation
  event subprocess starts should have first-class process-level semantics.
\item[Runtime fidelity] Whether to deepen true BPMN scope semantics as the
  project evolves.
\end{description}
